\begin{textsample}{POS Dim 1 – pe – Score 34.00 – pe/pe\_ef\_71.txt}  \label{ex:f1_pos_185}
6.4 HISTÓRIA O conhecimento histórico emerge como \textbf{disciplina} com pretensões científicas no \textbf{século} XIX.
Associada, em grande medida, ao aparecimento dos Estados \textbf{Modernos} e arrogando -se a responsabilidade de contar a história desses Estados e, por \textbf{conseguinte}, da formação da nação e do povo, necessários à constituição de cada \textbf{um} deles, ficava, portanto, reservada à História a função de estabelecer os marcos, os feitos, os fatos e o panteão de heróis \textbf{que} teriam sido responsáveis pela história da nação e da construção de \textbf{uma} dada identidade nacional.
Na análise de Elza Nadai : > O \textbf{século} XIX acrescentou, paralelamente aos grandes movimentos \textbf{que} ocorreram visando construir os Estados Nacionais sob hegemonia burguesa, a necessidade de retornar -se ao passado, com o objetivo de identificar a “ base comum ” formadora da nacionalidade.
Daí os conceitos tão caros às histórias nacionais : Nação, Pátria, Nacionalidade, Cidadania ( NADAI, 1986, p. 106 ).
Nesse sentido, a História tinha por finalidade forjar, mediante o estabelecimento de \textbf{uma} continuidade entre passado-presente-futuro, os traços da nacionalidade nos indivíduos de \textbf{um} dado território, proporcionando -lhes a construção de \textbf{uma} memória e de \textbf{uma} tradição.
Ela buscava assumir, assim, não só \textbf{uma} função científica na ordem do saber, mas também \textbf{um} sentido magisterial na formação dos \textbf{sujeitos} necessários aos \textbf{ideais} nacionalistas de composição dos Estados \textbf{Modernos}.
É nesse contexto \textbf{que} a História aparecerá, pela primeira vez, para o pensamento ocidental, como \textbf{uma} \textbf{disciplina} \textbf{que} precisa ser ensinada.
A História era pensada como mestra da vida.
Conhecer e aprender com a história tinha \textbf{um} claro sentido exemplar : fazer com \textbf{que} as gerações futuras não repetissem os \textbf{erros} do passado e aprendessem com os exemplos dos grandes \textbf{homens} em seus feitos heroicos e exemplares.
No entanto, essa modalidade de escrita e ensino de História trazia embutida em seus \textbf{fundamentos} \textbf{um} claro viés elitista de legitimação dos estratos superiores da sociedade, ligados ao establishment dos Estados \textbf{Modernos}, além de servir de narrativa de justificação de \textbf{uma} dada realidade política e social e sua reprodução.
Grosso modo, essa concepção de História se traduzia, reproduzia e circulava mediante \textbf{uma} prática de ensino positivista, tradicional, \textbf{que} ecoava \textbf{uma} história das e para as elites políticas, alçadas à condição de únicos \textbf{sujeitos} produtores da História.
Não tardou e no próprio \textbf{século} XIX ocorreram as primeiras críticas a essa modalidade de escrita da história e às finalidades ideológicas de sua reprodução.
Ainda naquele período, o filósofo alemão Karl Marx colocou -se como \textbf{um} crítico voraz daquela concepção e a acusou de estar a serviço de \textbf{uma} ideologia aristocrática, burguesa e excludente.
O marxismo propôs \textbf{uma} escrita da história politicamente engajada e comprometida com a conscientização da classe trabalhadora.
Nesse sentido, o ensino da História deveria servir à conscientização das massas para a revolução.
A história é investida não só de \textbf{um} sentido teleológico, mas também é pensada como orientadora da ação política para a transformação da realidade social e revolução da própria história, mediante o estabelecimento de \textbf{uma} nova ordem, de \textbf{um} novo tempo.
O ensino de história teria, em seu \textbf{horizonte}, a formação e a conscientização da classe trabalhadora, forjando \textbf{um} \textbf{sujeito} crítico e revolucionário \textbf{que} toma consciência de \textbf{que} é sujeito de sua própria história.
O marxismo retira das elites o privilégio de se \textbf{dizerem} os únicos fazedores da história.
Marx defende \textbf{que} “ os \textbf{homens} fazem a história, mesmo \textbf{que} não o saibam ”.
Seria função dos historiadores e/ou dos professores de história, como intelectuais de vanguarda, promover, mediante o processo de \textbf{ensino-aprendizagem}, essa tomada de consciência do \textbf{sujeito} histórico, fazendo -o crítico e autônomo para \textbf{que} pudesse agir e intervir politicamente sobre sua realidade e modificá -la ( REIS, 1999, 2003 ).
Mas, foi da Escola dos Annales, no início do \textbf{século} XX, \textbf{que} vieram as maiores alterações tanto na produção do conhecimento histórico e da sua escrita quanto nas implicações para o seu ensino.
É com os Annales \textbf{que} a História vai ser pensada, \textbf{talvez} pela primeira vez, em relação às demais ciências sociais.
Problematizada, de acordo com Marc Bloch, como “ a ciência dos \textbf{homens} no tempo ”, a História vai ser pensada na sua \textbf{interface} com a Antropologia, a Sociologia, a Psicologia, a Economia, a Ciência Política, etc.
Esse diálogo possibilita \textbf{uma} ampliação dos \textbf{métodos}, das técnicas, sobretudo dos problemas, dos objetos, das \textbf{abordagens} e, em especial, das fontes necessárias para sua escrita.
Isso amplia de forma considerável a dimensão e o quantitativo dos \textbf{sujeitos} \textbf{que} ganham \textbf{dignidade} historiográfica, ou seja, \textbf{que} passam a ser \textbf{vistos} não só como \textbf{sujeitos} históricos, mas também como produtores de narrativas sobre o passado.
Há, dessa maneira, \textbf{uma} quebra em relação aos modos de se produzir, conhecer e ensinar História estabelecidos, até então, sobretudo em relação à tradição dita positivista e sua concepção linear, teleológica e elitista de História ( REIS, 1999, 2003 ).
Essa tendência à pluralização e à diversificação dos estudos históricos se acentuou ainda mais a partir da segunda metade do \textbf{século} XX, com o desenvolvimento de \textbf{inúmeras} outras modalidades de produção e escrita da História.
Essas foram decorrentes da \textbf{crise} de paradigmas \textbf{que} \textbf{atravessou} as Ciências Humanas a partir daquele período e \textbf{que} questionava, radicalmente, as grandes narrativas \textbf{que} buscavam dar sentido à realidade e à própria história da \textbf{humanidade}.
A História não é mais \textbf{vista} como mestra da vida.
Ela se investe de outras funções cada vez mais específicas e contextualizadas a dados períodos e espaços.
Muito embora preserve \textbf{uma} função magisterial bem clara : “ fazer \textbf{lembrar} o \textbf{que} \textbf{uma} dada sociedade quer esquecer ”.
Dessa forma, a produção do conhecimento histórico vai ser diretamente relacionada aos problemas do presente.
O historiador é aquele \textbf{que} conhece o passado a partir dos problemas colocados no e pelo presente.
Ele passa a se pensar tanto como sujeito quanto como objeto do próprio conhecimento \textbf{que} produz em \textbf{uma} relação permanente de tensão entre presente e passado na qual este é constantemente reescrito à medida \textbf{que}, como defende Paul Veyne ( 1992 ), o repertório de questões muda ou se amplia.
É nesse sentido \textbf{que} o ensino de História passa a ser pensado como promotor das diferenças, problematizador de identidades \textbf{cristalizadas} e direcionado à promoção da cidadania, da pluralidade e da \textbf{democratização} dos modos de vida.
Diante dessa perspectiva, cabe destacar o pensamento de Schmidt e Cainelli : > É importante destacar \textbf{que} a História tem \textbf{uma} função didática de formar \textbf{uma} consciência histórica cada vez mais complexa, com a perspectiva de fornecer elementos para a orientação, interpretação do passado, para dentro, [ des]construindo identidades, e para fora, fornecendo sentidos para a ação na vida prática ( SCHIMIDT e CAINELLI, 2009, p. 37 ).
É, portanto, diante dessas perspectivas e contextos de produção do conhecimento histórico \textbf{que} os currículos de História, tanto em nível nacional quanto em níveis estaduais e municipais, começaram a ser pensados no Brasil, sobretudo a partir da década de 1980. \textbf{Ora} privilegiando \textbf{uma} dada \textbf{abordagem}, \textbf{ora} contemplando \textbf{uma} diversidade delas.
Prevista desde a LDBEN ( 1996 ), a construção de \textbf{uma} Base Nacional Comum Curricular foi objeto de discussão de educadores e especialistas ao longo das décadas seguintes, o \textbf{que} produziu \textbf{uma} série de documentos orientadores \textbf{que} vão dos Parâmetros Curriculares Nacionais, passando pelas Diretrizes Curriculares Nacionais até chegar às discussões \textbf{que} \textbf{fomentaram} a \textbf{atual} BNCC.
É nesse \textbf{bojo} de discussões \textbf{que} emerge a BNCC de História, trazendo em si todas estas tensões e disputas \textbf{inerentes} ao próprio campo do saber histórico, mas apontando para a necessidade de \textbf{um} curriculum mínimo \textbf{que} garanta o direito à aprendizagem em História de dados conteúdos, conceitos, práticas e o desenvolvimento de competências e habilidades \textbf{que} expressem a constituição de \textbf{um} dado sujeito, pensado como cidadão formador de \textbf{uma} sociedade \textbf{justa}, democrática e plural ( BNCC, 2017 ).
Esse desafio proposto pela BNCC coloca para os Estados alguns problemas, em especial para a construção do curriculum de História, \textbf{uma} vez \textbf{que} esses terão de compatibilizar as especificidades de suas realidades sociais, culturais e históricas com \textbf{um} conjunto de competências e habilidades gerais já estabelecidas pelo documento nacional a partir de \textbf{uma} \textbf{abordagem}, sobretudo nos anos finais do ensino fundamental, cronológica.
Apesar de estabelecer \textbf{que} esse é apenas \textbf{um} arranjo possível e \textbf{que} os Estados e municípios podem fazer ou construir outros, a BNCC dificulta a possibilidade de pensá -los a partir de eixos temáticos, como é feito nos “ Parâmetros Curriculares de História do Estado de Pernambuco ” ( PERNAMBUCO, 2013 ).
Dessa maneira, o currículo \textbf{que} se seguirá será pensado levando em consideração essas variadas \textbf{abordagens}, sem desconsiderar os entendimentos diversos, procurando deixar margem de flexibilidade para \textbf{que} os municípios e demais sistemas de ensino possam imprimir as características \textbf{que} julgarem necessárias quando este documento chegar ao chão da sala de aula.

% matched lemmas: abordagem, atravessar, atual, bojo, conseguinte, crise, cristalizar, democratização, dignidade, disciplina, dizer, ensino-aprendizagem, erro, fomentar, fundamento, homem, horizonte, humanidade, ideal, inerente, interface, inúmero, justo, lembrar, moderno, método, ora, que, sujeito, século, talvez, um, ver
\end{textsample}
